\section{Exact values at a prescribed initial spectrum}
\label{sec:spectralvalue94}
The equality set of a resource bound does not determine the value away
from that set.  We now solve this remaining optimization for the biased
basis experiment of Theorem~\ref{thm:operationalhierarchy90}.  Besides
giving the exact spectral loss, the solution distinguishes protocols
that merely retain a receiver from protocols that communicate a
receiver outcome to the next preparation.

Fix $d\ge2$, a weighted exact second-moment family satisfying
\eqref{eq:basismoments90}, and $0\le t\le1$.  The alternative basis is
selected once and used at both calls.  In this section only, the priors
are those of \eqref{eq:gameprior90}; write
\[
 L_0=2(d+1)-t^2,\qquad p_E=(d+1)/L_0.
\]
Fix the \emph{pure} first acquisition through its Gram matrix
$\rho=C_0^*C_0\in\mathcal D_d$.  Its input marginal is
$\rho^{\mathsf T}$.  The initial acquisition is not replaced by a
classically resolved mixture.  Subsequent instruments, mixtures of
fresh states, public seeds and tester limits retain the conventions of
Definition~\ref{def:rankclass92}.

For a nonnegative vector $u$, put $v_i=\sqrt{u_i}$ and define
\begin{equation}\label{eq:secularmatrix94}
 \chi(u)=\lambda_{\max}\bigl(vv^*-\diag u\bigr).
\end{equation}
Appending zeros does not change this nonnegative number.  In the secular
equation omit zero coordinates; the matrix definition remains valid on
every support face.  If at least
two entries of $u$ are positive, it is the unique positive solution of
\begin{equation}\label{eq:secularroot94}
                 \sum_i\frac{u_i}{u_i+\chi}=1.
\end{equation}
If at most one entry is positive, set $\chi(u)=0$.

Let $\lambda_1\ge\cdots\ge\lambda_d\ge0$ be the eigenvalues of
$\rho$.  For $1\le r\le d$, choose the smallest
$m\in\{0,\ldots,r-1\}$ for which
\[
 \lambda_{m+1}\le\tau,\qquad
 \tau=\frac{\sum_{i>m}\lambda_i}{r-m},
\]
and define the rank-$r$ majorant
\begin{equation}\label{eq:rankmajorant94}
 q^{(r)}(\lambda)
   =(\lambda_1,\ldots,\lambda_m,
               \underbrace{\tau,\ldots,\tau}_{r-m},0,\ldots,0).
\end{equation}
The choice always exists at $m=r-1$.  The preceding failed test, if
$m>0$, implies $\lambda_m>\tau$, so the displayed vector is decreasing.
Zero tails are allowed; ``rank-$r$'' here means rank at most $r$.
Write $\mathcal H_r(\rho)=\chi(q^{(r)}(\lambda(\rho)))$.
The ``smallest $m$'' convention fixes the active set at a water-level
tie.  Entries equal to $\tau$ belong to the flattened tail.  Including
such a tied entry in the initial segment would give the same vector,
so $q^{(r)}$ is unique even there.  If $r$ reaches the support size,
$q^{(r)}=\lambda$, padded by zeros.  These are the conventions for
$\chi$ and $q^{(r)}$ throughout the article; eigenbasis choices are not
part of their definitions.  Lemma~\ref{lem:spectralroof94} proves the
least-majorant assertion and hence uniqueness intrinsically.

\begin{lemma}[Continuity at water-level ties]\label{lem:watermap96}
On decreasing probability vectors the map $q^{(r)}$ is continuous,
piecewise affine, and nonexpansive in $\ell^1$:
\[
 \|q^{(r)}(\lambda)-q^{(r)}(\mu)\|_1
                                     \le\|\lambda-\mu\|_1.
\]
\end{lemma}
\begin{proof}
For each fixed active index $m$, the displayed rule is an affine map:
its first $m$ coordinates are unchanged, and each remaining input
coordinate contributes $1/(r-m)$ to each of the next $r-m$ outputs.
Its matrix has nonnegative entries and every column sums to one, so
its induced $\ell^1$ norm is one.  The finitely many closed active
regions are polyhedral.  On their overlaps the water-level equalities
make the formulas identical, including zero tails.  Hence the map is
continuous.  A line segment in the decreasing simplex has a finite
partition into these regions; summing the norm bounds along its
collinear subsegments proves the assertion.
\end{proof}

\begin{theorem}[Prescribed-spectrum value]\label{thm:spectralvalue94}
For the once-selected finite family and the fixed pure first acquisition
above, the optimum over the complete class $\mathfrak R_{k,\ell}$,
$1\le k,\ell\le d$, is
\begin{equation}\label{eq:spectralvalue94}
 P_{k,\ell}^{\rm b}(\rho)
       =p_E+\frac{t^2}{2L_0}\mathcal H_r(\rho),
                 \qquad r=\min\{k,\ell\}.
\end{equation}
At $r=1$ the spectral gain is zero.  At $t=0$ the value is $1/2$.
If $r\ge\operatorname{rank}\rho$, the majorant is the initial spectrum
itself.  Trace norms used in the derivation are unhalved; the factor
$1/2$ in \eqref{eq:spectralvalue94} converts the filtered-swap leaf
normalization into this Bayes gain.
The optimum is attained by an instrument with at most $d$ recorded
outcomes, the same instrument after every first device label, and fresh
references of dimension at most $r$.  Both retained quantum registers
may have dimension $r$.  The second preparation may depend on the
recorded receiver outcome, but need not depend on the first device
label.  This construction acts on the prescribed initial acquisition;
it does not optimize over its Gram matrix.
\end{theorem}
For $r=d$ no instrument is needed and the answer is given directly by
the secular equation for $\lambda(\rho)$.  The theorem gives an exact
value for the specified priors.  The equal-prior spectral rigidity of
Section~\ref{sec:spectralrigidity93} remains a separate result.

\subsection{The scalar function and a rank-constrained leaf}
\begin{lemma}[Secular concavity]\label{lem:secular94}
The function $\chi$ is continuous, symmetric, homogeneous of degree one,
coordinatewise nondecreasing and concave on each nonnegative orthant.
It is strictly Schur-concave on each fixed positive trace simplex:
if $u\prec w$ and $u,w$ are not permutations, then
$\chi(u)>\chi(w)$.  Here vectors are padded to a common length.
\end{lemma}
\begin{proof}
On a support containing at least two coordinates, the matrix in
\eqref{eq:secularmatrix94} has strictly positive off-diagonal entries.
A maximizing Rayleigh vector can be chosen nonnegative, and its
coordinates on this support must be positive.  Its eigenvalue equation
is $(u_i+\chi)x_i=\sqrt{u_i}\langle v,x\rangle$.
The Rayleigh value is positive, and summing gives
\eqref{eq:secularroot94}.  Its left side strictly decreases from the
support size to zero, proving uniqueness.  The support-one case follows
directly from the matrix formula.  That formula also gives continuity,
symmetry and homogeneity, including at the boundary.

For $a>0$, the root characterization says
\[
 \chi(u)\ge a\quad\Longleftrightarrow\quad
                     \sum_i\frac{u_i}{u_i+a}\ge1.
\]
Each summand is increasing and concave in $u_i$, so every such
superlevel set is convex and coordinatewise increasing.  If
$c=\chi(u)>0$ and $e=\chi(w)>0$, convexity of the unit superlevel set
applied with weights $c/(c+e),e/(c+e)$ gives
$\chi(u+w)\ge c+e$.  If one value vanishes the same inequality follows
from coordinate monotonicity.  Homogeneity now proves concavity.

Symmetry and concavity imply Schur concavity by convex combinations of
permutations.  For strictness, if $\chi(w)>0$, strict concavity of
$x/(x+\chi(w))$ shows that $u\prec w$, without permutation equality,
implies
$\sum_i u_i/(u_i+\chi(w))>1$.  The root for $u$ is therefore larger.
If $\chi(w)=0$, $w$ has only one positive entry and any nonpermutation
vector majorized by it has at least two, again giving strictness.
\end{proof}

For $A,B\succeq0$ let
\[
 \mathcal N(A,B)=\tr\bigl[-(\sqrt A\otimes\sqrt B)S
                                  (\sqrt A\otimes\sqrt B)\bigr]_+.
\]
\begin{lemma}[Optimal fresh Gram at a fixed old Gram]
\label{lem:spectralleaf94}
If $a_1\ge\cdots\ge a_d\ge0$ are the eigenvalues of $A$, then
\begin{equation}\label{eq:spectralleaf94}
 \max_{B\in\mathcal D_d,\ \operatorname{rank}B\le\ell}
                  \mathcal N(A,B)
           =\frac12\chi(a_1,\ldots,a_\ell).
\end{equation}
A maximizing $B$ commutes with $A$ and is supported on its leading
$\ell$ eigenspaces.  If $c=\chi(a_1,\ldots,a_\ell)>0$, an explicit
choice on that support is
\begin{equation}\label{eq:optimalfresh94}
 b_i=\frac{a_i/(a_i+c)^2}{\sum_{j\le\ell}a_j/(a_j+c)^2}.
\end{equation}
No invertibility of $A$ is required.  A commuting optimizer exists;
if a leading eigenspace is degenerate the chosen eigenbasis, and hence
the displayed optimizer, need not be unique.  The factor $1/2$ in
\eqref{eq:spectralleaf94} concerns negative mass, with no halving of
the trace norm in Lemma~\ref{lem:swapfilter90}.
\end{lemma}
\begin{proof}
Lemma~\ref{lem:swapfilter90} expresses $2\mathcal N(A,B)$ as
$(\tr\sqrt Z)^2-\tr Z$, where $Z=\sqrt A B\sqrt A$.
First work on the support of $A$.  Compressing $B$ to that support
leaves $Z$ unchanged and does not increase its rank.  If the compressed
trace is positive, normalizing it can only increase the nonnegative,
homogeneous objective.  If the trace vanishes the objective is zero.
We may therefore take $A$ invertible on its support and
$\tr(A^{-1}Z)=1$.

For fixed eigenvalues $z_1\ge\cdots\ge z_s>0$ of $Z$, $s\le\ell$,
rearrangement gives
\[
                  \sum_{i=1}^s z_i/a_i\le\tr(A^{-1}Z)=1.
\]
For completeness, in eigenbases the trace is a sum of products weighted
by squared entries of a unitary, a doubly stochastic matrix.  Its
minimum over that larger convex set is achieved at a permutation;
exchanging two inverted pairs puts the largest $z_i$ on the largest
$a_i$.  This proves the inequality.  Put $Z'$ diagonal with these
$z_i$ on the leading eigenspaces of $A$ and
$B'=A^{-1/2}Z'A^{-1/2}$.  Then $0<\tr B'\le1$;
normalizing $B'$ proves that a commuting maximizer with the asserted
support suffices.

For such a $B$, write $x_i=\sqrt{b_i}$, so $\|x\|_2=1$.  The objective
is
\[
 2\mathcal N(A,B)
       =\left(\sum_{i\le\ell}\sqrt{a_i}x_i\right)^2
                          -\sum_{i\le\ell}a_ix_i^2.
\]
The Rayleigh maximum is \eqref{eq:secularmatrix94}, with a nonnegative
maximizer.  Its eigenvector equation yields
\eqref{eq:optimalfresh94}.  With at most one positive $a_i$ the maximum
is zero, achieved by any rank-one Gram on the support, or any state if
$A=0$.  This handles every singular case explicitly.
\end{proof}

\subsection{A spectral evaluation of the concave envelope}
The next lemma records the majorization step separately.  It applies
to every continuous symmetric concave function, not only to the leaf
objective just derived.  The underlying majorization and ensemble
methods are classical; their operational use in entanglement
transformation and assistance is discussed in
\cite{Nielsen94,JonathanPlenio94,Assistance94}.

\begin{lemma}[Rank-capped spectral envelopes]\label{lem:spectralroof94}
Let $f$ be a continuous symmetric concave function on the probability
simplex of length $r$.  For a density matrix of rank at most $r$, pad
its eigenvalue vector to length $r$ before applying $f$.  Then
\begin{equation}\label{eq:spectralroof94}
 \max_{\substack{\rho=\sum_hw_ha_h\,,\ w_h\ge0\\
                         a_h\in\mathcal D_d^{\le r}}}
        \sum_hw_h f(\lambda(a_h))
        =f(q^{(r)}_1(\lambda(\rho)),\ldots,
                                   q^{(r)}_r(\lambda(\rho))).
\end{equation}
The maximum admits at most $d$ atoms commuting with $\rho$, all with
the same spectrum $q^{(r)}$.  In particular $\mathcal H_r$ is concave
on $\mathcal D_d$; its homogeneous extension to positive matrices is
concave and superadditive.
\end{lemma}
\begin{proof}
Write $q=q^{(r)}(\lambda)$.  Its construction gives $\lambda\prec q$.
Indeed its initial $m$ partial sums agree with those of $\lambda$;
each subsequent entry of $\lambda$ is at most $\tau$, and its total
sum is one.  Moreover $q$ is majorized by every decreasing probability
vector $z$ of rank at most $r$ which majorizes $\lambda$.  Up to $m$
this follows from the majorization inequalities.  Between $m$ and $r$
the decreasing entries of $z$ give
\[
 \sum_{i\le j}z_i\ge
 \frac{r-j}{r-m}\sum_{i\le m}z_i+\frac{j-m}{r-m}
 \ge\sum_{i\le m}\lambda_i+(j-m)\tau
                   =\sum_{i\le j}q_i.
\]
After $r$ both totals are one.  Thus $q$ is the least such majorant.

For any finite decomposition, the variational formula for the sum of
the $j$ largest eigenvalues gives
\[
 \lambda(\rho)\prec z:=\sum_hw_h\lambda(a_h),
\]
where the eigenvalue vectors are decreasing and padded to length $d$.
The vector $z$ has rank at most $r$.  Concavity and Schur concavity give
$\sum_hw_hf(\lambda(a_h))\le f(z)\le f(q)$.

For the converse, $\lambda\prec q$ implies that $\lambda$ lies in the
convex hull of the permutations of $q$.  One direct verification is
as follows.  For every real vector $b$, rearrangement followed by
summation by parts in the partial-sum inequalities gives
$b\cdot\lambda\le b^\downarrow\cdot\lambda
 \le b^\downarrow\cdot q=\max_\pi b\cdot\pi q$.
A separating hyperplane is therefore impossible.  This convex hull
lies in a hyperplane of dimension at most $d-1$, so Carath\'eodory
reduces the combination to at most $d$ permutations.  In an eigenbasis
of $\rho$ they give commuting density matrices $a_h$ of spectrum $q$.
Their $f$ values coincide, proving attainment.

Concavity follows also by combining any two attaining decompositions.
For a countable decomposition, the same upper follows from the Ky Fan
inequalities and Jensen's inequality by convergence of its bounded
trace weights.  At zero trace define the homogeneous extension to be
zero.  Trace-weighted convex combinations prove its concavity and
superadditivity.  Apply this with $f=\chi$ for the final assertion.
\end{proof}

For clarity, the homogeneous extension just proved is
\[
 \widehat{\mathcal H}_r(A)=
 \begin{cases}
  (\tr A)\mathcal H_r(A/\tr A),&\tr A>0,\\
  0,&A=0.
 \end{cases}
\]
For positive $A,B$ of traces $a,b>0$, concavity on densities gives
$\widehat{\mathcal H}_r(A+B)\ge
\widehat{\mathcal H}_r(A)+\widehat{\mathcal H}_r(B)$ by weights
$a/(a+b)$ and $b/(a+b)$; zero traces are immediate.  Subsequent formulas
write $\mathcal H_r(A)$ for this extension on positive matrices.

\subsection{Proof of the prescribed-spectrum theorem}
\begin{proof}[Proof of Theorem~\ref{thm:spectralvalue94}]
Use the dimension-preserving normal form of
Lemma~\ref{lem:ranknormal92}.  The fixed initial Gram gives
$\sum_hA_{yh}=\rho$ for every $y$.  Put $w_{yh}=\tr A_{yh}$ and
$a_{yh}=A_{yh}/w_{yh}$ on nonzero branches.  These are Gram weights;
they are not substituted for the hypothesis-dependent history
probabilities.  Mixed fresh states and receiver channels are refined
classically, without increasing either retained quantum dimension.

The biased payoff blocks are $-\kappa S$ for equal labels and
$-2\kappa\Pi_-/(d-1)\preceq0$ for unequal labels, where
$\kappa=t^2/(dL_0)$.  Thus only equal labels can improve on the
constant alternative decision.  Lemma~\ref{lem:spectralleaf94} gives
\[
 P-p_E\le\frac{\kappa}{2}\sum_{y,h}
                  \chi(\lambda_1(A_{yh}),\ldots,\lambda_\ell(A_{yh})).
\]
Each summand is at most $\mathcal H_r(A_{yh})$.  If $k\le\ell$, the
rank of $A_{yh}$ is at most $r=k$ and the two expressions agree.  If
$\ell<k$, the first $r=\ell$ coordinates of its rank-$r$ majorant
are coordinatewise at least its leading $r$ eigenvalues; use the
monotonicity of $\chi$.  Superadditivity then gives
\[
 P-p_E\le\frac{\kappa}{2}\sum_y\mathcal H_r(\rho)
                    =\frac{t^2}{2L_0}\mathcal H_r(\rho).
\]
The same bound survives public mixing, countable normal instruments
and tester closure, by the fixed-Gram normal form and trace tails.
At $t=0$ all devices coincide and both sides reduce to $1/2$.

For attainment take the commuting decomposition
$\rho=\sum_hw_ha_h$ supplied by Lemma~\ref{lem:spectralroof94}, all
$a_h$ with spectrum $q^{(r)}$.  On each support choose rectangular
$C_h,D_h:\C^d\to\C^r$ with
\[
 C_h^*C_h=w_ha_h,\qquad D_h^*D_h=b_h,
\]
where $b_h$ is the permuted optimizing fresh Gram in
\eqref{eq:optimalfresh94}.  If $\chi(q^{(r)})=0$, use any pure
fresh Gram and the constant decision.  Otherwise $b_h$ is a density
matrix of rank at most $r$, supported with $a_h$.  The maps
$V_h=C_hC_0^+$ satisfy $V_hC_0=C_h$ and
$\sum_hV_h^*V_h=I$ on the initial Schmidt support: positivity of the
Gram sum implies that every $C_h$ vanishes on $\ker\rho$.
They therefore form an instrument on the given initial receiver.
Proposition~\ref{prop:finiteinstrument96} gives its domains, codomains,
completion on an unused orthogonal complement and the at-most-$d$
message bound explicitly, without retaining an extra environment.

Use this same instrument after every $y$ and prepare $|D_h\rangle$
independently conditional on its recorded outcome.  For equal labels,
the positive spectral projection of the filtered negative swap attains
$\kappa w_h\chi(q^{(r)})/2$.  For unequal labels choose the zero
scalar-hypothesis effect.  Complements complete each binary readout.
Summing the $d$ first labels gives equality.  These unnormalized
formulas also define the endpoint $t=1$ and null histories.  The
construction proves attainment without an extra retained environment.
\end{proof}

\subsection{When a receiver message is strictly useful}
The following comparison fixes the initial acquisition and changes
only the permission to send a receiver outcome to the fresh source.
Let $P_{d,r}^{\rm nm}(\rho)$ be the optimum with old receiver dimension
$d$ and fresh reference dimension $r$, when the fresh state may depend
on the first device label and an independent public seed, but on no
outcome obtained by measuring the old receiver.  Arbitrary receiver
processing and the final joint readout remain allowed.  Any old
processing without such communication can be postponed to the final
readout by Lemma~\ref{lem:postponement96}; the old bound $d$ contains
the entire initial Schmidt support.  This is a subclass of $\mathfrak R_{d,r}$, not the
complete-measure-and-reprepare class.

\begin{corollary}[Exact value of a receiver message]\label{cor:messagevalue94}
For the fixed experiment and initial Gram above,
\begin{equation}\label{eq:nomessage94}
 P_{d,r}^{\rm nm}(\rho)
          =p_E+\frac{t^2}{2L_0}\chi(\lambda_1,\ldots,\lambda_r).
\end{equation}
If $t>0$, then
\begin{equation}\label{eq:messagestrict94}
 P_{d,r}^{\rm b}(\rho)>P_{d,r}^{\rm nm}(\rho)
                 \quad\Longleftrightarrow\quad
                      2\le r<\operatorname{rank}\rho.
\end{equation}
The strict gain is realized with at most $d$ receiver outcomes and
without first-device-label dependence of the fresh state.
\end{corollary}
\begin{proof}
Without a receiver message each first-label branch retains Gram
$\rho$.  Lemma~\ref{lem:spectralleaf94} gives the displayed upper,
and the commuting optimizer with the positive-part final decisions
attains it.  Public mixtures cannot improve an already maximized
linear score.  If $r=1$, both secular values vanish.  If
$r\ge\operatorname{rank}\rho$, no tail is moved and $q^{(r)}=\lambda$.
In the remaining case the leading $r$ eigenvalues are all positive,
and $q^{(r)}$ increases at least one of them without decreasing any.
For a vector with at least two positive entries the root strictly
increases when any positive coordinate is increased, as follows
immediately from \eqref{eq:secularroot94}.  This proves both directions.
\end{proof}

For an explicit example take $d=3$, $r=2$ and
$\rho=\diag(3/4,1/8,1/8)$.  Then
\[
 q^{(2)}=(3/4,1/4,0),\quad
 \chi(q^{(2)})=\sqrt3/4,\quad
 \chi(\lambda_1,\lambda_2)=\sqrt6/8.
\]
With $L_0=8-t^2$, the exact advantage of the message is
\begin{equation}\label{eq:messageexample94}
 \frac{t^2(2\sqrt3-\sqrt6)}{16L_0}>0\qquad(t>0).
\end{equation}
Use two Gram atoms
$a_1=\diag(3/4,1/4,0)$ and $a_2=\diag(3/4,0,1/4)$ with
weights $1/2$.  With $C_0=\sqrt\rho$, explicit old-receiver maps are
\[
 V_1=\begin{pmatrix}1/\sqrt2&0&0\\0&1&0\end{pmatrix},\qquad
 V_2=\begin{pmatrix}1/\sqrt2&0&0\\0&0&1\end{pmatrix}.
\]
They satisfy $\sum_hV_h^*V_h=I_3$ and
$C_h=V_hC_0$, $C_h^*C_h=a_h/2$.
The fresh Grams are respectively
$\diag(1/2,1/2,0)$ and $\diag(1/2,0,1/2)$; normalized fresh
vectors have the corresponding two equal Schmidt weights.  The final
positive-part readouts used in the proof attain the claimed scores.
Thus the improvement is a realizable recorded instrument, not a
change of initial state or postselection on a favorable history.

\begin{corollary}[Spectral loss and exact saturation]\label{cor:spectralsaturation94}
For $r>1$, the difference from the optimized rank-$r$ value of
Theorem~\ref{thm:rankspectrum92} is exactly
\begin{equation}\label{eq:exactspectralloss94}
 P_{k,\ell}^{\rm b}-P_{k,\ell}^{\rm b}(\rho)
       =\frac{t^2}{2L_0}\left(\frac{r-1}{r}-\mathcal H_r(\rho)\right).
\end{equation}
At $t>0$ it vanishes exactly when $\rho\preceq I/r$.
For $t>0$ and a fixed $\rho$ of rank $s$, its constrained optimum strictly
increases as the smaller receiver dimension increases from $1$ to
$s$, and is constant thereafter.  The smallest dimension attaining
its full-reset value is therefore $s$, including the trivial case
$s=1$.  At $t=0$ every dimension gives the same value $1/2$.
\end{corollary}
\begin{proof}
The uniform vector of length $r$ is majorized by every probability
vector of that length.  Lemma~\ref{lem:secular94} gives
$\chi(q)\le(r-1)/r$, with equality only at that uniform vector.
The construction of $q^{(r)}$ makes this equivalent to
$\lambda_1\le1/r$.  Subtracting \eqref{eq:spectralvalue94} from
\eqref{eq:rankbiased92} proves the exact loss identity.

The least-majorant property gives
$q^{(r+1)}\prec q^{(r)}$.  If $r<s$, these have respectively $r+1$
and $r$ positive entries and are not permutations, so strict Schur
concavity gives a strict increase.  At $r\ge s$ both vectors equal
$\lambda(\rho)$.  This proves the saturation statement.
\end{proof}

The secular root and the majorant are functions of a supplied spectrum.
Bisection of \eqref{eq:secularroot94} evaluates the value without an
optimization over physical protocols; a finite convex decomposition
then realizes it.  Real input spectra and real weights are permitted.
These statements do not assume a rational eigenbasis or give an
efficient construction of the unknown-device design.  In particular,
the general gain from a receiver message proved here is compatible
with the no-feedback globally attaining preparations in
Corollary~\ref{cor:threecuts93}: those preparations optimize the initial
Gram, whereas \eqref{eq:messagestrict94} holds it fixed.
